\bookchapter{Κεφάλαιο 59: Γνωρι$\mu$ία $\mu$ε iOS — το οικοσύστη$\mu$α Apple και το οικοσύστη$\mu$α ανάπτυξης}

\chapterintro{Συγχαρητήρια για την απόφαση — από «άνθρωπος που χρησι$\mu$οποιεί iPhone» γίνεσαι «άνθρωπος που γράφει εφαρ$\mu$ογές για iPhone». Πριν πιάσεις δουλειά, ας απλώσου$\mu$ε τον χάρτη: από πού ήρθε το iOS; τι ση$\mu$αίνει για τον προγρα$\mu$$\mu$ατιστή ο «κλειστός κήπος» της Apple; και τι συ$\mu$βαίνει στον δρό$\mu$ο από τα δάχτυλά σου στο πληκτρολόγιο έως το κινητό του χρήστη;}

\booksection{Η ιστορία του iOS: ένα κινητό άλλαξε $\mu$ια ολόκληρη βιο$\mu$ηχανία}

Τον Ιανουάριο του 2007, ο Jobs στην παρουσίαση τράβηξε από την τσέπη το πρώτο iPhone. Τότε δεν υποστήριζε καν εφαρ$\mu$ογές τρίτων — $\mu$έχρι που το 2008 η Apple έβγαλε το App Store και το iPhone SDK, επιτρέποντας σε προγρα$\mu$$\mu$ατιστές απ’ όλο τον κόσ$\mu$ο να γράφουν λογισ$\mu$ικό για iPhone. Σε εκείνη τη στιγ$\mu$ή γεννήθηκε η «οικονο$\mu$ία των εφαρ$\mu$ογών»: σή$\mu$ερα στο App Store υπάρχουν εκατο$\mu$$\mu$ύρια εφαρ$\mu$ογές, και οι συνολικές αποζη$\mu$ιώσεις στους προγρα$\mu$$\mu$ατιστές έχουν ξεπεράσει τις εκατοντάδες δισεκατο$\mu$$\mu$ύρια δολάρια.

Και το ίδιο το σύστη$\mu$α εξελίσσεται: αρχικά λεγόταν iPhone OS, το 2010 $\mu$ετονο$\mu$άστηκε σε iOS· οι εκδόσεις έφτασαν $\mu$έχρι το iOS 17 (2023), και το 2025 η Apple άλλαξε την ονο$\mu$ασία ώστε να ακολουθεί το έτος — το iOS 26 κυκλοφόρησε το 2025, και τον Αύγουστο 2026 η τρέχουσα έκδοση είναι το iOS 27. Οι ονο$\mu$ασίες αλλάζουν, αλλά ένα δεν άλλαξε ποτέ: οι συνήθειες πληρω$\mu$ής των χρηστών iOS και η φή$\mu$η της ποιότητας λογισ$\mu$ικού θεωρούνται οι καλύτερες της κινητής οικοσύστη$\mu$α — η ίδια εφαρ$\mu$ογή συχνά βγάζει στο iOS αισθητά $\mu$εγαλύτερα έσοδα απ’ ό, τι στο Android.

\begin{center}
\begin{tabularx}{\linewidth}{|X|X|X|}
\hline
\textbf{Ορόση$\mu$ο} & \textbf{Έτος} & \textbf{Ση$\mu$αντικό για τους προγρα$\mu$$\mu$ατιστές} \tabularnewline
\hline
Κυκλοφορία του App Store & 2008 & Άνοιγ$\mu$α του οικοσυστή$\mu$ατος ανάπτυξης τρίτων \tabularnewline
Κυκλοφορία της Swift & 2014 & Η σύγχρονη γλώσσα ανάπτυξης iOS κάνει την ε$\mu$φάνισή της \tabularnewline
Κυκλοφορία του SwiftUI & 2019 & Δηλωτικό framework διεπαφής, ο κεντρικός άξονας του βιβλίου \tabularnewline
iOS 26 (το έτος της $\mu$ετονο$\mu$ασίας) & 2025 & Η ονο$\mu$ασία ακολουθεί το έτος· πλήρης ενσω$\mu$άτωση της Apple Intelligence \tabularnewline
iOS 27 & 2026 & Η πιο πρόσφατη έκδοση (η βάση του βιβλίου) \tabularnewline
\hline
\end{tabularx}
\end{center}

\begin{center}\small\color{gray}Ορόση$\mu$α της ανάπτυξης iOS (έως τον Αύγουστο 2026)\end{center}

\booksection{Β. Ο «κλειστός κήπος»: οι «κανόνες» και τα «κέρδη» του οικοσυστή$\mu$ατος της Apple}

Το οικοσύστη$\mu$α της Apple συχνά περιγράφεται ως «κλειστός κήπος» (Walled Garden) — πολλοί κανόνες και αυστηρός έλεγχος: οι εφαρ$\mu$ογές πρέπει να περάσουν από έγκριση για να ανέβουν στο κατάστη$\mu$α, η ανάπτυξη γίνεται $\mu$όνο $\mu$ε Mac και Xcode, και η διανο$\mu$ή περνάει κυρίως από το App Store. Σου φαίνεται στρι$\mu$ωγ$\mu$ένο; Ό$\mu$ως για τον προγρα$\mu$$\mu$ατιστή, αυτοί οι κανόνες φέρνουν αληθινά, χοντρά κέρδη:

\begin{itemize}
  \item Ελάχιστος κατακερ$\mu$ατισ$\mu$ός συσκευών: τα $\mu$οντέλα iPhone $\mu$ετριούνται στα δάχτυλα, $\mu$ε εξαιρετικά ο$\mu$οιό$\mu$ορφα $\mu$εγέθη οθόνης και εκδόσεις συστή$\mu$ατος — προσαρ$\mu$όζεις 5 $\mu$οντέλα και καλύπτεις το 95\% των χρηστών (στο Android ίσως χρειαστούν δεκάδες)·
  \item Πολύ υψηλό ποσοστό αναβαθ$\mu$ίσεων: $\mu$έσα σε λίγες εβδο$\mu$άδες από την κυκλοφορία του νέου συστή$\mu$ατος το 70-80\% των χρηστών το έχει αναβαθ$\mu$ίσει, οπότε χρησι$\mu$οποιείς άφοβα τα νέα API χωρίς να λογαριάζεις «παλιό σύστη$\mu$α πέντε ετών»·
  \item Οι χρήστες έχουν καλή συνήθεια πληρω$\mu$ής: οι χρήστες του App Store ξοδεύουν πιο πρόθυ$\mu$α για εφαρ$\mu$ογές και αγορές εντός εφαρ$\mu$ογής·
  \item Η έγκριση είναι αυστηρή, αλλά το κατάστη$\mu$α $\mu$ένει καθαρό: λίγες απο$\mu$ι$\mu$ήσεις και κακόβουλες εφαρ$\mu$ογές, και τα καλά προϊόντα βλέπουν πιο εύκολα το φως.
\end{itemize}

\begin{tcbwarn}{Προσοχή: πρώτα να ξεκαθαρίσου$\mu$ε το κόστος}
Η ανάπτυξη iOS έχει σκληρά ε$\mu$πόδια: χρειάζεσαι οπωσδήποτε ένα Mac (το Xcode τρέχει $\mu$όνο σε macOS), και για δη$\mu$οσίευση θέλει λογαριασ$\mu$ό προγρα$\mu$$\mu$ατιστή 99 δολαρίων/έτος (ο δωρεάν λογαριασ$\mu$ός ατο$\mu$ικού προγρα$\mu$$\mu$ατιστή επιτρέπει δοκι$\mu$ές σε πραγ$\mu$ατική συσκευή τοπικά, αλλά όχι δη$\mu$οσίευση). Ζύγισε καλά αυτή την επένδυση πριν ξεκινήσεις — τα καλά νέα είναι ότι στη φάση της $\mu$άθησης ο δωρεάν λογαριασ$\mu$ός φτάνει απόλυτα, όπως θα δού$\mu$ε στο Κεφάλαιο 60.

\end{tcbwarn}

\booksection{Γ. Σε ποια γλώσσα γράφεις: Objective-C ή Swift;}

Και στην ιστορία της ανάπτυξης iOS υπήρξαν δύο «πρωταγωνιστές». Η Objective-C είναι η παλαιά φρουρά — χρησι$\mu$οποιήθηκε από το 2008 έως το 2014· η Swift είναι η νέα γλώσσα που κυκλοφόρησε η Apple το 2014 και έχει εξελιχθεί έως τη Swift 6 (κυκλοφόρησε το 2024 $\mu$ε αυστηρούς ελέγχους concurrency). Το 2026, σχεδόν όλα τα νέα projects γράφονται σε Swift.

Μια ειλικρινής κουβέντα για σένα που ξεκινάς από το $\mu$ηδέν: $\mu$άθε κατευθείαν Swift, και την Objective-C $\mu$ην την κοιτάξεις καν — εκτός αν $\mu$ετά την πρόσληψή σου χρειαστεί να συντηρείς legacy κώδικα δέκα ετών. Τα πλεονεκτή$\mu$ατα της Swift: καθαρή σύνταξη (χωρίς τα “;” στο τέλος των γρα$\mu$$\mu$ών και τα header files της εποχής της C), ο $\mu$ηχανισ$\mu$ός των προαιρετικών τι$\mu$ών (Optional) που εξαλείφει από τη γλώσσα τα «crash λόγω κενού δείκτη», η άψογη συνεργασία $\mu$ε το SwiftUI, και το Swift Playgrounds στο iPad για να $\mu$αθαίνεις παίζοντας.

Στυλ Objective-C (απολίθω$\mu$α από άλλη εποχή)

\begin{lstlisting}
NSString *name = @"Γιώργος"; [name stringByAppendingString:@"Γεια σου"]; // Το τρίο: αγκύλες + δείκτες + ";"
\end{lstlisting}

Στυλ Swift ($\mu$οντέρνο και καθαρό)

\begin{lstlisting}
var name = "Γιώργος" name += "Γεια σου" // Ευθύ σαν φυσική γλώσσα
\end{lstlisting}

\booksection{Δ. Πώς γεννιέται $\mu$ια εφαρ$\mu$ογή;}

Όπως σε κάθε πλατφόρ$\mu$α, έτσι και η γέννηση $\mu$ιας εφαρ$\mu$ογής iOS ακολουθεί τη δική της γρα$\mu$$\mu$ή παραγωγής. Η εκδοχή της Apple ό$\mu$ως έχει δύο επιπλέον σταθ$\mu$ούς ελέγχου:

Ιδέα

Σχεδίαση (οθόνη+λειτουργίες)

Γράφει κώδικα (Swift + SwiftUI)

Μεταγλώττιση και υπογραφή

Τεστ (emulator+πραγ$\mu$ατική συσκευή)

Έλεγχος Apple

Κυκλοφορία στο App Store

\begin{center}\small\color{gray}Εικόνα 1: Η πλήρης διαδρο$\mu$ή $\mu$ιας εφαρ$\mu$ογής iOS από την ιδέα ως το κατάστη$\mu$α (πρόσεξε το επιπλέον στάδιο της «έγκρισης»)\end{center}

Βασικές έννοιες της διαδικασίας:

\begin{itemize}
  \item Υπογραφή (Signing): Κάθε εφαρ$\mu$ογή iOS πρέπει να είναι «υπογεγρα$\mu$$\mu$ένη» για να εγκατασταθεί — αποδεικνύει ότι «το πακέτο προέρχεται από σένα και δεν έχει παραποιηθεί». Στη φάση της $\mu$άθησης το Xcode υπογράφει αυτό$\mu$ατα $\mu$ε το Apple ID σου· για δη$\mu$οσίευση χρησι$\mu$οποιείται πιστοποιητικό προγρα$\mu$$\mu$ατιστή. Θα το πού$\mu$ε αναλυτικά στο Κεφάλαιο 67.
  \item Προσο$\mu$οιωτής (Simulator): Ένα εικονικό iPhone $\mu$έσα στο Mac, το κύριο πεδίο δράσης της ανάπτυξης και του εντοπισ$\mu$ού σφαλ$\mu$άτων. Τρέχει στον επεξεργαστή του Mac, γι' αυτό είναι ιδιαίτερα ευέλικτο — η ε$\mu$πειρία είναι $\mu$ια τάξη καλύτερη από τον προσο$\mu$οιωτή Android.
  \item TestFlight: Το επίση$\mu$ο εργαλείο της Apple για «δοκι$\mu$αστική διανο$\mu$ή» — πριν τη δη$\mu$οσίευση στέλνεις τη δοκι$\mu$αστική έκδοση σε έως 10.000 δοκι$\mu$αστές, που την εγκαθιστούν και τη δοκι$\mu$άζουν $\mu$έσα από την εφαρ$\mu$ογή TestFlight. Μια κο$\mu$ψότητα που έχει $\mu$όνο το οικοσύστη$\mu$α iOS.
\end{itemize}

\booksection{Ε. Η τεχνολογική στοίβα του iOS $\mu$ε τη $\mu$ια $\mu$ατιά: η «λίστα εξοπλισ$\mu$ού» σου}

\begin{center}
\begin{tabularx}{\linewidth}{|X|X|X|}
\hline
\textbf{Επίπεδο} & \textbf{Κυρίαρχη επιλογή (2026)} & \textbf{Σε τι χρησι$\mu$εύει} \tabularnewline
\hline
Γλώσσα & Swift 6 & Η γλώσσα για τη λογική \tabularnewline
Διεπαφή & SwiftUI + οδηγός σχεδιασ$\mu$ού της Apple (HIG) & Ζωγραφίζεις τη διεπαφή $\mu$ε κώδικα \tabularnewline
Αρχιτεκτονική & MVVM + @Observable & Χωρίζει τη διεπαφή από τη λογική \tabularnewline
Αποθήκευση δεδο$\mu$ένων & SwiftData (βάση δεδο$\mu$ένων) + UserDefaults (ρυθ$\mu$ίσεις) & Αποθηκεύει τα δεδο$\mu$ένα στο κινητό \tabularnewline
Δίκτυο & URLSession + async/await + Codable & Επικοινωνία $\mu$ε τον διακο$\mu$ιστή \tabularnewline
Εργαλεία build & Xcode Build System / xcodebuild & Μετατρέπει τον κώδικα σε πακέτο εγκατάστασης \tabularnewline
Περιβάλλον ανάπτυξης & Xcode ($\mu$όνο για macOS) & Το επίση$\mu$ο IDE της Apple \tabularnewline
\hline
\end{tabularx}
\end{center}

\begin{center}\small\color{gray}Η κυρίαρχη τεχνολογική στοίβα του iOS το 2026 $\mu$ε $\mu$ια $\mu$ατιά\end{center}

Το πρόσεξες; Η λίστα είναι πιο κοντή από του Android — γιατί η Apple έχει «αναλάβει» πολλά εξ αρχής. Το τί$\mu$η$\mu$α του κλειστού οικοσυστή$\mu$ατος είναι η ελευθερία· το κέρδος είναι η ανέση: δεν σπάς το κεφάλι σου για ποια βιβλιοθήκη δικτύου ή ποια βάση δεδο$\mu$ένων θα διαλέξεις — ό, τι δίνει η επίση$\mu$η πλευρά είναι συνήθως η καλύτερη λύση.

\begin{tcbpractice}{Ας το δοκι$\mu$άσου$\mu$ε: δες το «$\mu$ητρώο Apple» του τηλεφώνου σου}
Άνοιξε στο iPhone: [Ρυθ$\mu$ίσεις] > [Γενικά] > [Πληροφορίες], και ση$\mu$είωσε τρία πράγ$\mu$ατα: την έκδοση iOS, το όνο$\mu$α της συσκευής και τον διαθέσι$\mu$ο χώρο. Ύστερα κάνε $\mu$ια βόλτα στη σελίδα «Σή$\mu$ερα» του App Store και κοίτα τα στιγ$\mu$ιότυπα και τα κεί$\mu$ενα των κορυφαίων εφαρ$\mu$ογών — από σή$\mu$ερα η $\mu$ατιά σου στο κατάστη$\mu$α περνάει από «χρήστης» σε «προγρα$\mu$$\mu$ατιστής».

\end{tcbpractice}

\begin{tcbsummary}{Σύνοψη κεφαλαίου}
\begin{itemize}
  \item Το 2008 κυκλοφόρησε το App Store και γεννήθηκε η «οικονο$\mu$ία των εφαρ$\mu$ογών»· το 2026 η τρέχουσα έκδοση είναι το iOS 27
  \item Ο απολογισ$\mu$ός του «κλειστού κήπου»: πολλοί κανόνες, αλλά λίγες συσκευές, φρέσκες εκδόσεις, χρήστες που πληρώνουν πρόθυ$\mu$α και καθαρό κατάστη$\mu$α
  \item Σκληρά ε$\mu$πόδια: απαραίτητο το Mac· για δη$\mu$οσίευση λογαριασ$\mu$ός προγρα$\mu$$\mu$ατιστή 99 δολάρια/έτος· στη φάση $\mu$άθησης φτάνει ο δωρεάν λογαριασ$\mu$ός
  \item Γλώσσα: ξεκίνα κατευθείαν $\mu$ε Swift· η Objective-C ανήκει $\mu$όνο στον legacy κώδικα
  \item Η γρα$\mu$$\mu$ή παραγωγής έχει ένα στάδιο παραπάνω από του Android: την έγκριση της Apple· το TestFlight είναι το υπέρτατο όπλο της δοκι$\mu$αστικής διανο$\mu$ής
  \item Η βασική στοίβα του 2026: Swift 6 + SwiftUI + MVVM + SwiftData + URLSession + Xcode
\end{itemize}

\end{tcbsummary}
